% Generated from locale/tr/content/sets-functions-relations/sets/proofs-about-sets.tex; do not hand-edit.


\olfileid[tr]{sfr}{set}{prf}
\olsection{Kümeler Hakkında İspatlar}

\begin{editorial}
Bu bölümün yerini \verb|content/methods/proofs| almıştır ve bölüm varsayılan
olarak bu kitaptan çıkarılmıştır.
\end{editorial}

\begin{explain}
Kümeler ve şimdiye kadar tanıttığımız gösterimler, kümeleri belirtmek ve
aralarındaki bağıntıları ifade etmek için elverişli kısaltmalar sağlar. Çoğu
zaman bu bağıntılar hakkındaki iddiaları ispatlamak da gerekir. Matematiksel
ispatlara aşina değilseniz bu sizin için yeni olabilir. Bu nedenle basit bir
örneği adım adım inceleyeceğiz. Her $X$ ve $Y$ kümesi için her zaman
$X \cap (X \cup Y) = X$ olduğunu ispatlayacağız. Kümeler arasındaki böyle bir
özdeşlik nasıl ispatlanır? Kümelerin yalnızca !!{element}s{}ıyla belirlendiğini,
yani kümelerin ancak ve ancak aynı !!{element}s{}ı taşıdıklarında özdeş olduğunu
hatırlayın. Bu durumda (a) $X \cap (X \cup Y)$ kümesinin her !!{element}{}ının
$X$'in de !!{element}{}ı olduğunu ve tersine (b) $X$'in her !!{element}{}ının
$X \cap (X \cup Y)$ kümesinin de !!{element}{}ı olduğunu göstermeliyiz. Başka
bir deyişle hem (a) $X \cap (X \cup Y) \subseteq X$ hem de
(b) $X \subseteq X \cap (X \cup Y)$ olduğunu gösteririz.

``$X \cap (X \cup Y)$ kümesinin her~$z$ !!{element}{}ı, $X$'in de
!!{element}{}ıdır'' gibi genel bir iddia, önce gelişigüzel bir
$z \in X \cap (X \cup Y)$ verildiği varsayılarak ve bu varsayımdan
$z \in X$ sonucu çıkarılarak ispatlanır. Bu kalıbı ``genel koşullu ispat''
adıyla biliyor olabilirsiniz. Bu ispata katılan, örneğin ``$\cap$'' ve
``$\cup$'' gibi kavramların tanımlarını da kullanmamız gerekecek. Dolayısıyla
(a) durumu şöyle gerekçelendirilir:

\begin{quote}
(a) Önce $X \cap (X \cup Y) \subseteq X$ olduğunu, yani $\subseteq$ tanımı
uyarınca her~$z$ için $z \in X \cap (X \cup Y)$ ise $z \in X$ olduğunu
göstermek istiyoruz. $z \in X \cap (X \cup Y)$ olduğunu varsayın. $z$, iki
kümenin kesişiminin !!{element}{}ı ancak ve ancak her iki kümenin de
!!{element}{}ı olduğundan $z \in X$ ve ayrıca $z \in X \cup Y$ sonucuna
varabiliriz. Özellikle $z \in X$; göstermek istediğimiz de buydu.
\end{quote}

Bu, ispatın ilk yarısını tamamlar. Son adımda bir varsayımdan bir tümel
evetleme ($z \in X$ ve $z \in X \cup Y$) çıkıyorsa, her bileşenin de aynı
varsayımdan çıktığı olgusunu kullandığımıza dikkat edin. Bu kuralı ``tümel
evetleme giderimi'' ya da $\land$Elim adıyla biliyor olabilirsiniz. Şimdi de
(b)'yi ispatlayalım:

\begin{quote}
(b) Şimdi $X \subseteq X \cap (X \cup Y)$ olduğunu, yani $\subseteq$ tanımı
uyarınca her~$z$ için $z \in X$ ise $z \in X \cap (X \cup Y)$ olduğunu
ispatlayalım. $z \in X$ varsayımını yapın. $z \in X \cap (X \cup Y)$
olduğunu göstermek için ``$\cap$'' tanımı gereği (i) $z \in X$ ve ayrıca
(ii) $z \in X \cup Y$ olduğunu göstermeliyiz. Burada (i) yalnızca
varsayımımızdır; dolayısıyla başka bir şey ispatlamamız gerekmez. (ii) için
$z$'nin kümeler birleşiminin !!{element}{}ı ancak ve ancak bu kümelerden en az
birinin !!{element}{}ı olduğunu hatırlayın. $z \in X$ ve $X \cup Y$, $X$ ile
$Y$'nin birleşimi olduğundan burada durum budur. O halde $z \in X \cup Y$.
Hem (i) $z \in X$ hem de (ii) $z \in X \cup Y$ olduğunu gösterdik;
dolayısıyla ``$\cap$'' tanımı gereği $z \in X \cap (X \cup Y)$.
\end{quote}

Bu anlatım biraz uzundu, fakat kümeler ve aralarındaki bağıntılar hakkında
nasıl akıl yürüttüğümüzü gösterir. Genellikle bu kadar açık yazmayız;
özellikle bütün tanımları tekrarlamayabiliriz. Daha ileri bir metindeki ispat
çok daha sıkıştırılmış olur ve aşağıdakine benzeyebilir.
\end{explain}


\begin{prop}[Yutma]
Bütün $X$, $Y$ kümeleri için
\[
X \cap (X \cup Y) = X.
\]
\end{prop}

\begin{proof}
(a) $z \in X \cap (X \cup Y)$ olsun. O halde $z \in X$; dolayısıyla
$X \cap (X \cup Y) \subseteq X$.

(b) Şimdi $z \in X$ olsun. O zaman $z \in X \cup Y$ ve dolayısıyla
$z \in X \cap (X \cup Y)$. Böylece $X \subseteq X \cap (X \cup Y)$.
\end{proof}

\begin{prob}
$X \cup (X \cap Y) = X$ olduğunu ayrıntılı biçimde ispatlayın. Ardından
kısaltılmış, sıkıştırılmış bir ispat verin. (İpucu: $X \cup (X \cap Y)
\subseteq X$ yönü için durumlara ayırarak ispatı, diğer adıyla $\lor$Elim
kuralını kullanmanız gerekecek.)
\end{prob}

